\documentclass[aspectratio=169,10pt]{beamer}
\usetheme[numbering=fraction,progressbar=frametitle]{metropolis}
\usepackage{fontspec}
\setsansfont{Latin Modern Sans}
\usepackage{amsmath,amssymb,mathtools,graphicx,booktabs}
\usepackage{polyglossia}
\setdefaultlanguage{french}
\setbeamertemplate{navigation symbols}{}
\setbeamerfont{frametitle}{size=\large}
\setbeamerfont{normal text}{size=\normalsize}
\setlength{\parskip}{.45em}
\hypersetup{pdfauthor={},pdfsubject={Lecture slides},pdfkeywords={},pdftitle={Logique}}
\providecommand{\N}{\mathbb{N}}
\providecommand{\Z}{\mathbb{Z}}
\providecommand{\D}{\mathbb{D}}
\providecommand{\Q}{\mathbb{Q}}
\providecommand{\R}{\mathbb{R}}
\providecommand{\C}{\mathbb{C}}
\providecommand{\vect}[1]{\overrightarrow{#1}}
\providecommand{\vi}{\vec{\imath}}
\providecommand{\vj}{\vec{\jmath}}
\providecommand{\vk}{\vec{k}}
\providecommand{\norm}[1]{\left\lVert #1\right\rVert}
\providecommand{\vc}[2]{\begin{pmatrix}#1\\#2\end{pmatrix}}
\providecommand{\deter}[2]{\det\left(#1,#2\right)}
\providecommand{\abs}[1]{\left\lvert #1\right\rvert}
\providecommand{\dd}{\,\mathrm{d}}
\providecommand{\e}{\mathrm{e}}
\providecommand{\Card}{\operatorname{Card}}
\providecommand{\E}{\operatorname{E}}
\providecommand{\V}{\operatorname{V}}
\providecommand{\degre}{^{\circ}}
\providecommand{\un}[1]{\,\mathrm{#1}}
\providecommand{\sci}[2]{#1\times10^{#2}}
\providecommand{\pgcd}{\operatorname{PGCD}}
\providecommand{\ppcm}{\operatorname{PPCM}}
\title{Logique}
\subtitle{Cours et exemples}
\date{}
\author{}
\begin{document}
\begin{frame}[plain]
\vfill
{\usebeamerfont{title}\inserttitle\par}
\medskip
{\usebeamerfont{subtitle}\insertsubtitle\par}
\medskip
{\usebeamercolor[fg]{progress bar}\rule{\textwidth}{.5pt}}
\vfill
\end{frame}
\begin{frame}{}
Par groupes, classer dix phrases proposées par le professeur. Pour chaque phrase
classée vraie ou fausse, rédiger une justification; pour chaque phrase refusée,
expliquer quelle information manque. Transformer ensuite deux phrases ouvertes
en propositions à l'aide de « pour tout » ou « il existe ».
\end{frame}
\begin{frame}{Objectifs}
\begin{itemize}

  \item Distinguer énoncé, proposition et prédicat.
  
\item Maîtriser les connecteurs logiques et leurs tables de vérité.
  
\item Utiliser les quantificateurs \(\forall\) et \(\exists\).
  
\item Former la négation et la réciproque d'une implication.
  
\item Découvrir, en approfondissement, la contraposée et le raisonnement par l'absurde.

\end{itemize}
\end{frame}
\begin{frame}{Énoncés et propositions}
Une \textbf{proposition} est un énoncé pouvant être \textbf{vrai} (\textbf{V})
ou \textbf{faux} (\textbf{F}), mais pas les deux simultanément.
\end{frame}
\begin{frame}{Énoncés et propositions}
\begin{itemize}

  \item  \(2 + 2 = 4\)  : proposition vraie.
  
\item  \(7\) est pair  : proposition fausse.
  
\item  \(x > 3\)  : prédicat (sa valeur dépend de \(x\)), \emph{pas} une proposition.
  
\item   Viens ici  : phrase impérative, pas une proposition.

\end{itemize}
\end{frame}
\begin{frame}{Négation et connecteurs logiques}
Soient \(P\) et \(Q\) deux propositions. On définit :
\end{frame}
\begin{frame}{Négation et connecteurs logiques (2)}
\begin{center}\begin{tabular}{llllll}
\(P\) & \(Q\) & \(\lnot P\) & \(P \land Q\) & \(P \lor Q\) & \(P \Rightarrow Q\) \\ 
\textbf{V} & \textbf{V} & \textbf{F} & \textbf{V} & \textbf{V} & \textbf{V} \\ 
\textbf{V} & \textbf{F} & \textbf{F} & \textbf{F} & \textbf{V} & \textbf{F} \\ 
\textbf{F} & \textbf{V} & \textbf{V} & \textbf{F} & \textbf{V} & \textbf{V} \\ 
\textbf{F} & \textbf{F} & \textbf{V} & \textbf{F} & \textbf{F} & \textbf{V} \\ 
\end{tabular}\end{center}
\end{frame}
\begin{frame}{Négation et connecteurs logiques (3)}
L'implication \(P \Rightarrow Q\) est fausse \emph{uniquement} quand \(P\) est vraie et
\(Q\) est fausse. Le   ou  mathématique est \textbf{inclusif}.
\end{frame}
\begin{frame}{Négation et connecteurs logiques}
\(P \Leftrightarrow Q\) est vraie si \(P\) et \(Q\) ont la même valeur de vérité :
\[P \Leftrightarrow Q \;\equiv\; (P \Rightarrow Q) \land (Q \Rightarrow P).\]
\end{frame}
\begin{frame}{Négation et connecteurs logiques}
La \textbf{contraposée} de \(P \Rightarrow Q\) est \(\lnot Q \Rightarrow \lnot P\).
Ces deux implications sont toujours équivalentes. Pour démontrer
\(P \Rightarrow Q\), il suffit de démontrer \(\lnot Q \Rightarrow \lnot P\).
\end{frame}
\begin{frame}{Quantificateurs}
\begin{itemize}

  \item \(\forall x \in E,\ P(x)\) :   pour tout \(x\) dans \(E\), \(P(x)\) est vraie. 
  
\item \(\exists x \in E,\ P(x)\) :   il existe au moins un \(x\) dans \(E\) tel que
        \(P(x)\) est vraie. 

\end{itemize}
\end{frame}
\begin{frame}{Quantificateurs}
\[\lnot\bigl(\forall x \in E,\ P(x)\bigr) \;\equiv\; \exists x \in E,\ \lnot P(x)\]
\[\lnot\bigl(\exists x \in E,\ P(x)\bigr) \;\equiv\; \forall x \in E,\ \lnot P(x)\]
\end{frame}
\begin{frame}{Types de raisonnement}
Pour démontrer \(P \Rightarrow Q\) directement :
on suppose \(P\) vraie et on en déduit, par une chaîne d'implications,
que \(Q\) est vraie.
\end{frame}
\begin{frame}{Types de raisonnement}
On démontre \(\lnot Q \Rightarrow \lnot P\) à la place de \(P \Rightarrow Q\).
\textbf{Étapes :} (1) Supposer \(\lnot Q\). (2) En déduire \(\lnot P\).
\end{frame}
\begin{frame}{Types de raisonnement}
Pour démontrer \(P\) :
on suppose \(\lnot P\) et on aboutit à une \textbf{contradiction}.
La contradiction montre que \(\lnot P\) est fausse, donc \(P\) est vraie.
\end{frame}
\begin{frame}{Types de raisonnement}
Démontrons que \(\sqrt{2}\) est irrationnel.
\emph{Supposons} \(\sqrt{2} \in \mathbb{Q}\) : il existe \(p, q \in \mathbb{N}^*\) avec \(\operatorname{pgcd}(p,q)=1\)
et \(\sqrt{2} = p/q\).
Alors \(2q^2 = p^2\), donc \(p^2\) est pair, donc \(p\) est pair : \(p = 2k\).
On obtient \(2q^2 = 4k^2\), d'où \(q^2 = 2k^2\), et \(q\) est pair.
Mais alors \(2 \mid \operatorname{pgcd}(p,q)\), ce qui contredit \(\operatorname{pgcd}(p,q)=1\). \textbf{Contradiction.}
Donc \(\sqrt{2} \notin \mathbb{Q}\).
\end{frame}
\end{document}
